% P10-Core-Inhaltsprüfung, 15.09.2026: datierter Ergebnisanschluss; B-60 korrigiert.
% Belege: thesis-evidence/w2/core-content-audit-20260915/evaluation/
% P10-HITL, 15.09.2026: Entscheidungszweck und belegte Zustandswirkung präzisiert.
% P10-4b, 14.09.2026: Begriffs-/Beleganschlüsse und abgegrenzte Prüfaussagen präzisiert.
% P10-4a, 14.09.2026: Mechanismus, Quellenreichweite und v4-Nachweis angeglichen.
% M19 (12.09.2026): Leserführung und präzise Ergebnisdarstellung;
% Messwerte und historische Nachweisgrenzen bleiben erhalten.
% ============================================================
% §7.5 Integrierte Projektfortschreibung und Nachvollziehbarkeit
% v01 (Claude, 07.09.) — NEUE 8er-Struktur (Schreibplan-Zeile 7.5:
% Fallserie + Provenienzaudit + Anschlussmatrix v2).
% Grundlage: Z2-Fallblätter rev2 (nach Kollegen-Review: F5a/F5b
% getrennt, F1 feld-verglichen, F3 aufgeteilt, F4 n.b.).
% Budget ~2 S. + 1 Tabelle. Leitplanken: keine Erfolgsquote ·
% Falltypen ≠ unabhängige Läufe · Quer-Befunde sind Ergebnisse.
% ============================================================
% HERKUNFT: Fallblätter = thesis-evidence/w2/z2-fallblaetter.md
% (rev2; Läufe: 20260818_084712_765eea [F1/F5a], 20260817_123342_
% f23f06 [F2], 20260804_121433_eb181a [F3], 20260818_123055_a0b172
% [F4], 20260805_131538_899463 + Kommentar-Rundweg-Serie 20.08.
% [F5]; F1-Feldvergleich: 07-ingest/plan.json + state/core REQ-82 +
% github-snapshot 20260818_125520_cb157c Issues #44/#45);
% Provenienzaudit = runs/e2e-evidenz/traceability-audit.json
% (Skript v3, T1–T9; Core-Snapshot 06.09., 237 Items);
% Anschlussanalyse = kapitel-7/05-vergleichs-inspektion/
% 07-ANSCHLUSSFAEHIGKEIT.md (v2, qualitativ).
% ------------------------------------------------------------

\section{Integrierte Projektfortschreibung und Nachvollziehbarkeit}
\label{sec:eval-fortschreibung}

Ob die Kette den fachlichen Projektbestand \emph{richtig fortschreibt}, lässt
sich nicht allein an Abdeckungszahlen ablesen. Den Ausgangspunkt
bilden fünf gezielt gewählte fachliche Falltypen an
historisch belegten Läufen --- eine diagnostische Fallserie mit
retrospektiv dokumentierten Prüfkriterien, keine repräsentative
Stichprobe; eine Erfolgsquote wird bewusst nicht gebildet. Berichtet
werden fünf fachliche Falltypen und ein gesonderter externer
Rückwegfall (daher sechs Tabellenzeilen); zwei Fälle teilen
denselben Lauf --- es handelt sich nicht um zusätzliche
unabhängige Stichproben. Je Fall
wurden Einordnung, Inhaltserhaltung, Freigabebezug, Herkunft und
Projektionsentsprechung einzeln beurteilt
(Tabelle~\ref{t:eval-fallserie}); fehlende Belege führen zu
„nicht beurteilbar``, bewusst begrenzte Prüfumfänge zu „nicht
untersucht`` --- keiner dieser Zustände zählt als Erfüllung.
Für F1 und F3 wurde der Freigabebezug am 15.09.2026 ergänzend
anhand der gespeicherten Vorlagen, Entscheidungen, Anwendungsberichte
und Core-Zwischenstände abgeglichen. Die Nachprüfung verwendet die
historischen Artefakte; sie ist keine neue Systemausführung.

% P10-Core-Inhaltsprüfung, 15.09.2026: datierter Ergebnisanschluss; B-60 korrigiert.
% Belege: thesis-evidence/w2/core-content-audit-20260915/evaluation/
% B-58, 15.09.2026: 201 -> 207 -> 208 -> 212 Items, keine Löschung im geprüften Abschnitt.
% Nachtrag: thesis-evidence/w2/f1-zaehlkorrektur-20260915.md; historische Runs unverändert.
% ============================================================
% Tabelle 7.5-1: Ergebnisübersicht der Fallserie — v01 (07.09.)
% In Overleaf ablegen als: tables/Chap7-v02/chap7.5table1.tex
% Aufruf im Fließtext: % P10-Core-Inhaltsprüfung, 15.09.2026: datierter Ergebnisanschluss; B-60 korrigiert.
% Belege: thesis-evidence/w2/core-content-audit-20260915/evaluation/
% B-58, 15.09.2026: 201 -> 207 -> 208 -> 212 Items, keine Löschung im geprüften Abschnitt.
% Nachtrag: thesis-evidence/w2/f1-zaehlkorrektur-20260915.md; historische Runs unverändert.
% ============================================================
% Tabelle 7.5-1: Ergebnisübersicht der Fallserie — v01 (07.09.)
% In Overleaf ablegen als: tables/Chap7-v02/chap7.5table1.tex
% Aufruf im Fließtext: % P10-Core-Inhaltsprüfung, 15.09.2026: datierter Ergebnisanschluss; B-60 korrigiert.
% Belege: thesis-evidence/w2/core-content-audit-20260915/evaluation/
% B-58, 15.09.2026: 201 -> 207 -> 208 -> 212 Items, keine Löschung im geprüften Abschnitt.
% Nachtrag: thesis-evidence/w2/f1-zaehlkorrektur-20260915.md; historische Runs unverändert.
% ============================================================
% Tabelle 7.5-1: Ergebnisübersicht der Fallserie — v01 (07.09.)
% In Overleaf ablegen als: tables/Chap7-v02/chap7.5table1.tex
% Aufruf im Fließtext: \input{tables/Chap7-v02/chap7.5table1}
% Label: t:eval-fallserie
% Stil: chap6.1-Konvention.
% HERKUNFT: thesis-evidence/w2/z2-fallblaetter.md rev2
% (Ergebnistabelle; Urteils-Legende dort im Kopf). Läufe je Zeile
% siehe chap7.5.tex-HERKUNFT.
% ============================================================

\begin{table}[htbp]
  \centering
  \small
  \begin{tabular}{>{\raggedright\arraybackslash}p{0.16\textwidth}>{\raggedright\arraybackslash}p{0.12\textwidth}llll}
    \hline
    \textbf{Falltyp} & \textbf{Lauf (Stand)} & \textbf{Inhalt} & \textbf{Freigabe} & \textbf{Herkunft} & \textbf{Projektion} \\
    \hline
    Neue Anforderung & 18.08. & erfüllt\textsuperscript{a} & erfüllt\textsuperscript{a} & erfüllt & erfüllt (real) \\
    Präzisierung & 17.08. & teilweise\textsuperscript{f} & erfüllt & erfüllt & Dry-Run \\
    Widerspruch & 04.08. & erfüllt\textsuperscript{b} & erfüllt\textsuperscript{c} & \textbf{abweichend} & nicht untersucht \\
    Wiederholung & 18.08. & n.\,b. & erfüllt & erfüllt & n.\,a. \\
    Offene Frage & 18.08. & erfüllt & erfüllt & erfüllt & n.\,a.\textsuperscript{d} \\
    Ext.\ Rückweg & 20.08. & erfüllt & erfüllt & erfüllt & Hauptkette erfüllt\textsuperscript{e} \\
    \hline
    \multicolumn{6}{p{0.95\textwidth}}{\footnotesize
      Die Einordnung wurde fallbezogen beurteilt; bei F2 ist die
      Rücknahme des Kalender-Prüfauftrags nicht gesondert begründet
      (Fußnote~f). \textsuperscript{a}~für die verfolgten
      Ketten/belegten Operationen, kein Pauschalurteil über alle
      11 vor der GitHub-Projektion neu hinzugekommenen Core-Items.
\textsuperscript{b}~Ablösung und Relationen;
      fachliche Angleichung der abhängigen Backlog-Items blieb
      offener Klärungsbedarf. \textsuperscript{c}~nur das
      Entscheidungs-Gate lief interaktiv, die übrigen vier Gates im
      deklarierten Experiment-Modus. \textsuperscript{d}~offene
      Entscheidungen haben vertragsgemäß keinen Projektionskanal.
      \textsuperscript{e}~ein Vermerk-Schreibvorgang zunächst
      fehlgeschlagen; Run-Zuordnung des Nachtrags n.\,b.
      \textsuperscript{f}~Nachprüfung vom 15.09.: neue Vorgaben
      erhalten, alter Kalender-Prüfauftrag nur historisiert;
      Rücknahme nicht gesondert begründet. Dry-Run = Projektion
      vorbereitet, nicht ausgeführt.
      Legende: n.\,b. = Beleg fehlt · n.\,a. = vertragsgemäß nicht
      vorgesehen · „nicht untersucht`` = Prüfumfang bewusst
      begrenzt --- keine dieser Angaben ist eine Erfüllung.} \\
    \hline
  \end{tabular}
  \caption[Ergebnisse der fachlichen Fallserie]{Ergebnisübersicht der fachlichen Fallserie (fünf
    Falltypen, historisch belegte Läufe mit je eigenem
    Systemstand; keine Erfolgsquote). Beobachtete Abweichungen und
    Nachweislücken bleiben getrennt ausgewiesen: die
    Herkunfts-\emph{Anzeige} im Widerspruchsfall (abweichend), der
    zunächst fehlgeschlagene Vermerk-Schreibvorgang im Rückweg-Fall
    (Fußnote~e) sowie die n.\,b.-Zellen.}
  \label{t:eval-fallserie}
\end{table}

% Label: t:eval-fallserie
% Stil: chap6.1-Konvention.
% HERKUNFT: thesis-evidence/w2/z2-fallblaetter.md rev2
% (Ergebnistabelle; Urteils-Legende dort im Kopf). Läufe je Zeile
% siehe chap7.5.tex-HERKUNFT.
% ============================================================

\begin{table}[htbp]
  \centering
  \small
  \begin{tabular}{>{\raggedright\arraybackslash}p{0.16\textwidth}>{\raggedright\arraybackslash}p{0.12\textwidth}llll}
    \hline
    \textbf{Falltyp} & \textbf{Lauf (Stand)} & \textbf{Inhalt} & \textbf{Freigabe} & \textbf{Herkunft} & \textbf{Projektion} \\
    \hline
    Neue Anforderung & 18.08. & erfüllt\textsuperscript{a} & erfüllt\textsuperscript{a} & erfüllt & erfüllt (real) \\
    Präzisierung & 17.08. & teilweise\textsuperscript{f} & erfüllt & erfüllt & Dry-Run \\
    Widerspruch & 04.08. & erfüllt\textsuperscript{b} & erfüllt\textsuperscript{c} & \textbf{abweichend} & nicht untersucht \\
    Wiederholung & 18.08. & n.\,b. & erfüllt & erfüllt & n.\,a. \\
    Offene Frage & 18.08. & erfüllt & erfüllt & erfüllt & n.\,a.\textsuperscript{d} \\
    Ext.\ Rückweg & 20.08. & erfüllt & erfüllt & erfüllt & Hauptkette erfüllt\textsuperscript{e} \\
    \hline
    \multicolumn{6}{p{0.95\textwidth}}{\footnotesize
      Die Einordnung wurde fallbezogen beurteilt; bei F2 ist die
      Rücknahme des Kalender-Prüfauftrags nicht gesondert begründet
      (Fußnote~f). \textsuperscript{a}~für die verfolgten
      Ketten/belegten Operationen, kein Pauschalurteil über alle
      11 vor der GitHub-Projektion neu hinzugekommenen Core-Items.
\textsuperscript{b}~Ablösung und Relationen;
      fachliche Angleichung der abhängigen Backlog-Items blieb
      offener Klärungsbedarf. \textsuperscript{c}~nur das
      Entscheidungs-Gate lief interaktiv, die übrigen vier Gates im
      deklarierten Experiment-Modus. \textsuperscript{d}~offene
      Entscheidungen haben vertragsgemäß keinen Projektionskanal.
      \textsuperscript{e}~ein Vermerk-Schreibvorgang zunächst
      fehlgeschlagen; Run-Zuordnung des Nachtrags n.\,b.
      \textsuperscript{f}~Nachprüfung vom 15.09.: neue Vorgaben
      erhalten, alter Kalender-Prüfauftrag nur historisiert;
      Rücknahme nicht gesondert begründet. Dry-Run = Projektion
      vorbereitet, nicht ausgeführt.
      Legende: n.\,b. = Beleg fehlt · n.\,a. = vertragsgemäß nicht
      vorgesehen · „nicht untersucht`` = Prüfumfang bewusst
      begrenzt --- keine dieser Angaben ist eine Erfüllung.} \\
    \hline
  \end{tabular}
  \caption[Ergebnisse der fachlichen Fallserie]{Ergebnisübersicht der fachlichen Fallserie (fünf
    Falltypen, historisch belegte Läufe mit je eigenem
    Systemstand; keine Erfolgsquote). Beobachtete Abweichungen und
    Nachweislücken bleiben getrennt ausgewiesen: die
    Herkunfts-\emph{Anzeige} im Widerspruchsfall (abweichend), der
    zunächst fehlgeschlagene Vermerk-Schreibvorgang im Rückweg-Fall
    (Fußnote~e) sowie die n.\,b.-Zellen.}
  \label{t:eval-fallserie}
\end{table}

% Label: t:eval-fallserie
% Stil: chap6.1-Konvention.
% HERKUNFT: thesis-evidence/w2/z2-fallblaetter.md rev2
% (Ergebnistabelle; Urteils-Legende dort im Kopf). Läufe je Zeile
% siehe chap7.5.tex-HERKUNFT.
% ============================================================

\begin{table}[htbp]
  \centering
  \small
  \begin{tabular}{>{\raggedright\arraybackslash}p{0.16\textwidth}>{\raggedright\arraybackslash}p{0.12\textwidth}llll}
    \hline
    \textbf{Falltyp} & \textbf{Lauf (Stand)} & \textbf{Inhalt} & \textbf{Freigabe} & \textbf{Herkunft} & \textbf{Projektion} \\
    \hline
    Neue Anforderung & 18.08. & erfüllt\textsuperscript{a} & erfüllt\textsuperscript{a} & erfüllt & erfüllt (real) \\
    Präzisierung & 17.08. & teilweise\textsuperscript{f} & erfüllt & erfüllt & Dry-Run \\
    Widerspruch & 04.08. & erfüllt\textsuperscript{b} & erfüllt\textsuperscript{c} & \textbf{abweichend} & nicht untersucht \\
    Wiederholung & 18.08. & n.\,b. & erfüllt & erfüllt & n.\,a. \\
    Offene Frage & 18.08. & erfüllt & erfüllt & erfüllt & n.\,a.\textsuperscript{d} \\
    Ext.\ Rückweg & 20.08. & erfüllt & erfüllt & erfüllt & Hauptkette erfüllt\textsuperscript{e} \\
    \hline
    \multicolumn{6}{p{0.95\textwidth}}{\footnotesize
      Die Einordnung wurde fallbezogen beurteilt; bei F2 ist die
      Rücknahme des Kalender-Prüfauftrags nicht gesondert begründet
      (Fußnote~f). \textsuperscript{a}~für die verfolgten
      Ketten/belegten Operationen, kein Pauschalurteil über alle
      11 vor der GitHub-Projektion neu hinzugekommenen Core-Items.
\textsuperscript{b}~Ablösung und Relationen;
      fachliche Angleichung der abhängigen Backlog-Items blieb
      offener Klärungsbedarf. \textsuperscript{c}~nur das
      Entscheidungs-Gate lief interaktiv, die übrigen vier Gates im
      deklarierten Experiment-Modus. \textsuperscript{d}~offene
      Entscheidungen haben vertragsgemäß keinen Projektionskanal.
      \textsuperscript{e}~ein Vermerk-Schreibvorgang zunächst
      fehlgeschlagen; Run-Zuordnung des Nachtrags n.\,b.
      \textsuperscript{f}~Nachprüfung vom 15.09.: neue Vorgaben
      erhalten, alter Kalender-Prüfauftrag nur historisiert;
      Rücknahme nicht gesondert begründet. Dry-Run = Projektion
      vorbereitet, nicht ausgeführt.
      Legende: n.\,b. = Beleg fehlt · n.\,a. = vertragsgemäß nicht
      vorgesehen · „nicht untersucht`` = Prüfumfang bewusst
      begrenzt --- keine dieser Angaben ist eine Erfüllung.} \\
    \hline
  \end{tabular}
  \caption[Ergebnisse der fachlichen Fallserie]{Ergebnisübersicht der fachlichen Fallserie (fünf
    Falltypen, historisch belegte Läufe mit je eigenem
    Systemstand; keine Erfolgsquote). Beobachtete Abweichungen und
    Nachweislücken bleiben getrennt ausgewiesen: die
    Herkunfts-\emph{Anzeige} im Widerspruchsfall (abweichend), der
    zunächst fehlgeschlagene Vermerk-Schreibvorgang im Rückweg-Fall
    (Fußnote~e) sowie die n.\,b.-Zellen.}
  \label{t:eval-fallserie}
\end{table}


\textbf{Neue Anforderung (F1).} Ein gezielt gebautes,
synthetisches Meeting-Transkript durchlief die volle Kette bis zu
drei real erzeugten GitHub-Issues: Issue~\#43 betrifft die
Tagesübersicht dokumentierter Medikamentengaben; Issues~\#44 und~\#45
gehören zu dem in Abschnitt~\ref{sec:endtoend-pfade}
durchgearbeiteten Besuchsbeispiel. Die verfolgte Aussage
--- Besuchsankündigung „vorab \ldots{} mit Datum und Uhrzeit`` ---
blieb inhaltlich vom maschinellen Vorschlag über das autorisierte
Core-Item bis in den Issue-Titel erhalten. Der vollständige
Anforderungstext stimmt zwischen Vorschlag und REQ-82 wortgleich
überein. Von zwölf maschinellen Operationen wurden acht menschlich
übernommen und vier mit protokollierter Begründung abgewiesen.
Die Kennungen der acht freigegebenen Eingangsoperationen stimmen
genau mit den acht angewandten Operationen im Bericht überein;
die vier abgewiesenen kommen dort nicht vor. Ihr Vorlagetext und
ihre Begründungen sind als Ablehnungsvermerke im Core gespeichert.
Damit bleibt auch nachvollziehbar, was nicht übernommen wurde
und warum.

% P10-6: F1: menschliche Ablehnung und beobachtete Wirkung nicht durch die Tabelle unterbrechen.
\begin{samepage}
Der Operationsplan hatte die formale Prüfung bestanden. Dennoch
wies der Mensch die vorgeschlagene Verfeinerung von REQ-42 als
„Reine Bestätigung des Bestands`` ab: Der Ersatztext hätte die
vorhandenen Angaben zu einnahmebezogenen Bemerkungen verdrängt.
Dies betrifft den geltenden Anforderungstext; die bisherige
Fassung wäre bei einer Verfeinerung historisiert worden. Der
Vergleich der tatsächlich gespeicherten Core-Zwischenstände zeigt
dagegen unveränderten Text und unveränderte Version von REQ-42.
Damit ist die Wirkung der menschlichen Auswahl an diesem Fall
nachvollziehbar; die Verteilung acht/vier ist keine Qualitätsquote.
Grenze: Der Umschalter, der
den realen Schreibvorgang im wiederaufgenommenen Lauf aktivierte,
hinterlässt kein eigenes Artefakt --- der Anwendungsbericht belegt den
Schreibvorgang, nicht dessen Konfigurationsweg.
\par
\end{samepage}

\textbf{Präzisierung (F2).} Ein Autor-Diktat konkretisierte eine
vage Bestandsanforderung. Statt eines Duplikats entstand eine neue
Fassung von REQ-42 mit erhaltener Historie. Sie enthält die diktierten
Vorgaben: Sieben-Tage-Übersicht der Medikamenteneinnahmen für freigegebene
Angehörige zur eigenen Bezugsperson, ausschließlich Status und Zeitpunkt, keine Dosierungsdetails
und nur lesender Zugriff. Die ergänzende Inhaltsprüfung vom 15.09.2026
begrenzt jedoch das Erhaltungsurteil: Der frühere Prüfauftrag zur
Kalenderfunktion steht anschließend nur noch in der Historie; seine
Rücknahme ist nicht gesondert begründet. Titel, Zieltext und
Akzeptanzkriterien des abhängigen Backlog-Items standen bereits
wortgleich im Agentenvorschlag.
Belegt ist eine vorgeschlagene und menschlich freigegebene Angleichung,
keine eigenständige menschliche Textänderung. Die GitHub-Projektion
wurde nur als Dry-Run vorbereitet, nicht ausgeführt
(Anhang~\ref{anh:core-inhaltspruefung}, Kontrolle K02).

\textbf{Widerspruch (F3).} Der Konflikt betraf die Geltung von
No-Go-Einträgen, also Hinweisen auf zu vermeidende Themen oder
Handlungen. Die bestehende Anforderung REQ-32 verlangte eine
bewohnerbezogene Führung, die neue Aussage eine globale Geltung.
Die maschinelle Einordnung erfasste dies als offene Entscheidung.
Am interaktiven Entscheidungs-Gate übernahm der Mensch die neue
Festlegung und formulierte den endgültigen Text. Dieser steht
wortgleich in der neuen Anforderung REQ-78. Im nachfolgenden
Core-Zwischenstand ist REQ-32 als abgelöst erhalten, DEC-001
aufgelöst und die Deckungsbeziehung der PBIs PBI-012 und PBI-023
auf REQ-78 umgestellt. Der menschliche Freigabenachweis betrifft
hier ausschließlich das Entscheidungs-Gate; die übrigen vier
Gates des historischen Versuchs nutzten automatische Annahmen
im deklarierten Experimentbetrieb.
Das Ergebnis ist geteilt zu berichten: Erkennung,
Entscheid und Relationswirkung sind belegt; die fachliche
Angleichung der abhängigen Backlog-Items blieb als markierter
Klärungsbedarf offen, die Projektion wurde in diesem Experiment
nicht untersucht (Dry-Run). Zwei Befunde bleiben stehen: Die
Herkunfts-\emph{Anzeige} am Entscheidungspunkt nannte einen
fremden Lauf (Vorbefüllungs-Artefakt), und der berichtete
Core-Restore des Experiments hat kein maschinelles Artefakt.

\textbf{Wiederholte Information (F4).} Die Zweit-Einspeisung
desselben Deltas erzeugte keine neuen fachlichen Items und keine
Duplikate --- der Bestand wurde als solcher erkannt und die
zugehörige Entscheidung nicht dupliziert. Ob auch die
Bestandstexte unverändert blieben, wurde nicht per
Vorher-Nachher-Vergleich geprüft (n.\,b.); geprüft ist die
Delta-, nicht die Transkript-Ebene.

\textbf{Offene Frage und externer Rückweg (F5).} Eine
Meeting-Frage wurde als offene Entscheidung in den Bestand
aufgenommen und blieb offen --- es entstand kein erfundener
Beschluss; offene Fragen besitzen vertragsgemäß keinen eigenen
Projektionskanal. Auf dem Rückweg wurde ein externer
GitHub-Kommentar abgerufen, als Entwurf am Gate vorgelegt,
menschlich übernommen und bis zum aktualisierten Issue
zurückprojiziert; ein zweiter Kommentar wurde begründet abgewiesen
und beim Wiederauftauchen automatisch als „schon einmal
abgelehnt`` erkannt. Belegt ist der Kommentar-Eingang
(plus eine Konventions-Erkennung und ein menschlich aufgelöster
Drift-Fall an einem Issue) --- keine allgemeine
Änderungsüberwachung; ein Vermerk-Schreibvorgang schlug zunächst
still fehl und gelang erst im dritten Durchlauf.

\textbf{Übergreifende Befunde.} Die Fälle zeigen auch Grenzen der
Nachvollziehbarkeit: Die Rücksetzungen des Core nach Experimenten
sind nur in Notizen beschrieben; die angezeigte Herkunft kann von
der tatsächlichen abweichen; der Ausführungs-Umschalter hinterlässt
kein Artefakt. Offene Entscheidungen haben vertragsgemäß keine
Außensicht. Die Wiederholungsprüfung betrifft ausschließlich Deltas.

\subsection{Provenienzaudit}
\label{sec:eval-provenienz}

Ein deterministisches Audit untersucht die Herkunft der 237 Items
eines historischen Core-Snapshots. Nach der ersten Auswertung wurde
das Instrument am 13.09.2026 überarbeitet: Fassung~v4 unterscheidet
direkte Referenzen, exakt zugeordnete Eingangsobjekte und gerichtete
Beiträge zu solchen Zielen. Sie nutzt ausdrücklich gespeicherte
Lauf- und Dateibeziehungen; ein globaler Treffer derselben Kennung
genügt nicht. Der Systembestand blieb unverändert. Es handelt sich
um eine nachträgliche Verfeinerung der Auswertung, nicht um eine
Leistungsverbesserung des Systems.

% P10-4a, 14.09.2026: Audit v4; unterschiedliche Bezugsgrößen nicht addieren.
% Overleaf: tables/Chap7-v02/chap7.5table2.tex
\begin{table}[htbp]
  \centering
  \small
  \begin{tabular}{>{\raggedright\arraybackslash}p{0.45\textwidth}>{\raggedright\arraybackslash}p{0.43\textwidth}}
    \hline
    \textbf{Prüfgegenstand} & \textbf{Ergebnis und Reichweite} \\
    \hline
    Mindestens ein direkter Definitions- oder Eingangsbezug & 147/237 Items; keine Vollständigkeit aller Bezüge \\
    Direkter Bezug oder gerichteter Beitragspfad & 234/237 Items; indirekte Beiträge können Teilinhalte erklären \\
    Genau zugeordnetes Eingangsobjekt & 39 Items, davon 37 mit Plan-/Apply-Zuordnung \\
    Gespeicherte Ledger-Herkunftskanten & 104/104 im ausdrücklichen Ursprungskontext auflösbar \\
    Aktuelle Einzelreferenzvorkommen & 168/232 aufgelöst; 64 offen, verteilt auf 33 Items \\
    \hline
  \end{tabular}
  \caption[Getrennte Endpunkte des Herkunftsaudits v4]{Getrennte
    Endpunkte des Herkunftsaudits v4. Item-Mengen überschneiden sich;
    Kanten und Referenzvorkommen sind andere Zähleinheiten.
    Technische Auflösung ist kein semantischer Stützungsnachweis.}
  \label{t:provenienz-v4}
\end{table}


Tabelle~\ref{t:provenienz-v4} trennt die Bezugsgrößen. Ein
Beitragspfad kann etwa ein PBI zu einer Anforderung mit Quellenbezug
führen und dabei nur einen Teil seines Inhalts erklären.
Die 234 erreichten Items sind deshalb keine 234 vollständig
nachgewiesenen aktuellen Inhalte. Bei den drei übrigen Items ist
unter den engeren Prüfregeln nur ein eigenes Erzeugungsartefakt
auffindbar. Solche Dokumente zählen nicht als vorgelagerte Quelle.
Unter den identifizierten Eingängen finden sich Autor-, GitHub-
und Analyst-Eingaben; ihre Auflösung erfordert keine künstlichen
Ledger-Claims. Die Gruppen überschneiden sich.

% P10-6: Auditbefund samt semantischer Reichweitengrenze zusammenhalten.
\begin{samepage}
Alle 104 gespeicherten Ledger-Herkunftskanten sind im expliziten
Ursprungskontext auflösbar. Daneben bleiben 64 aktuelle
Einzelreferenzvorkommen bei 33 Items offen. Ein anderer erfolgreicher
Pfad hebt diese offenen Referenzen nicht auf. Die Quellenprüfung
unterscheidet zudem technische Auflösung und Textübereinstimmung:
Von 173 Zitaten der 71 eindeutig aufgelösten Claim-Definitionen sind
52 vollständig normalisiert im gebundenen Transkript auffindbar;
bei 104 passt nur der Textkörper nach dem frühen Doppelpunkt.
Die übrigen Fälle und die getrennte Prüfung der Unit-Bezüge stehen
in Anhang~\ref{anh:provenienz-v4}. In diesem Audit erfolgte keine neue semantische
Annotation.
\par
\end{samepage}

Auch der jüngste gespeicherte Übergang der 38 Items mit Version
größer eins wurde gesondert geprüft. Für 22 liegt ein Plan-/Apply-Beleg
mit übereinstimmendem primärem Textfeld vor, für zwei PBIs nur eine
Teilübereinstimmung. Bei 14 fehlt der Anschluss innerhalb der
unterstützten Regeln; neun davon besitzen unveränderten Haupttext.
Das ist weder ein Voll-Audit sämtlicher Änderungen noch der Nachweis
von 14 verlorenen Inhaltsquellen. Das Verfahren wurde mit 34
künstlichen Kontrollen geprüft und zweimal byteidentisch ausgeführt.
Anhang~\ref{anh:provenienz-v4} dokumentiert Verfahren, Belege und
Abgrenzung zur früheren Auswertung.

\subsection{Inhaltliche Prüfung ausgewählter Fortschreibungen}
\label{sec:eval-core-inhalte}

Auflösbare Referenzen belegen noch nicht, dass die verbundenen
Inhalte fachlich zusammenpassen. Ergänzend wurden deshalb am
15.09.2026 acht historische Einordnungsoperationen aus vier
Eingangsgruppen geprüft: Transkript, Autor-Delta, GitHub und sonstiges
Delta. Die Auswahlregel und die Kriterien wurden
vor ihrer Bewertung festgehalten; vier bereits bekannte Fälle
dienten getrennt als Kontrollen. Verglichen wurden Eingang,
Agentenvorschlag, Entscheidung, gespeicherter Inhalt und Beziehungen.
Anhang~\ref{anh:core-inhaltspruefung} dokumentiert Auswahl,
Einzelbefunde und Belegzugang.

Ein zusätzlicher Fall führt das Besuchsbeispiel weiter (N07).
Ein Agent leitete aus dem vorhandenen Projektwissen eine neue
Entwurfsregel ab: Angehörige sollen angefragte oder bestätigte
Besuche bis zu deren Beginn ändern oder absagen können. Änderungen
setzen den Status auf „angefragt`` zurück; nach einer Absage entfällt
die Erinnerung. Der Vorschlag ist als Ableitung gekennzeichnet,
nicht als Aussage aus dem Meeting. Nach menschlicher Freigabe
entstand REQ-91. Die nachfolgende, ebenfalls freigegebene
Arbeitspaketfortschreibung ergänzte das bestehende PBI-046 um diese
Regeln und erhielt dessen bisherige Vorgaben zur Ankündigung mit
Datum und Uhrzeit. Gespeicherte Beziehungen ordnen die neue
Anforderung diesem Arbeitspaket und dem Besuchs-Feature zu.
Der Projektionsbericht dokumentiert die Aktualisierung von
Issue~\#44. Damit sind Herkunft, fachliche Erweiterung und passende
Verbindungen an derselben Kette belegt.

Ein weiterer Fall zeigt die Wirkung einer Nichtübernahme (N03).
Das Autor-Delta verlangte für Besuchsankündigungen genau die drei
Statuswerte „angefragt``, „bestätigt`` und „abgelehnt``. Sie fehlten
im Agentenvorschlag. Der Mensch benannte diese Auslassung in der
Ablehnung; der fachliche Bestand blieb unverändert. Diese
unvollständige Vorlage wurde somit nicht wirksam. Dagegen wurde
in N06 eine mehrteilige GitHub-Frage im Agentenvorschlag einer
vorhandenen Entscheidung zugeordnet, die nur eine Teilfrage abdeckte. Die anschließende
Ablehnung belegt dort keine fachliche Korrektur; ihr Hintergrund
bleibt unsicher. Die Fallprüfung unterscheidet deshalb passende
Zuordnung, autorisierte Wirkung und fachliche Richtigkeit. Sie
liefert keine allgemeine Erfolgsquote für Agenten oder menschliche
Entscheidungen.

\subsection{Anschlussfähigkeit der Übergaben}
\label{sec:eval-anschluss}

Eine qualitative Abhängigkeitsanalyse dokumentiert, welche
Verträge die nachgelagerten Stufen von ihrem Eingang tatsächlich
benötigen --- im Kern Routinginformation (Disposition),
Quellbezug (Unit-Referenzen) und konventionskonforme Claim-Kennungen
innerhalb des jeweiligen Laufbestands. Die freie
Vergleichsbedingung liefert Aussagen mit Unit-Referenzen und
lauflokalen IDs; für die unmittelbare Übergabe an die bestehenden
Abnehmer fehlen insbesondere die vorgesehenen Dispositionen und
Facetten, und die Identitäten entsprechen nicht unmittelbar den
verwendeten Kanonisierungs- und Provenienzkonventionen. Das ist
eine Vertrags-, keine Fähigkeitsaussage und belegt keine
Alternativlosigkeit der gewählten Implementierung. Die
vollständige Feldzuordnung steht in
Tabelle~\ref{t:anh-vertragsfelder}
(Anhang~\ref{anh:evaluation}).
